\begin{textsample}{POS Dim 3 – 2009\_c4 – Score 1.00 – 2009\_c4/enade\_2009\_tec\_des\_mod\_3927bb5ec2.txt}  \label{ex:f3_pos_128}
Capítulo 4 Percepção sobre a Prova As percepções dos estudantes sobre a prova aplicada no ENADE/2009 na área de TECNOLOGIA EM DESIGN DE MODA foram mensuradas por meio de nove itens que avaliaram desde o grau de dificuldade da prova até o tempo gasto para concluí -la.
As questões foram analisadas separando -se concluintes e ingressantes, e as percepções sobre a prova foram relacionadas com o desempenho dos estudantes e com a Região Geográfica de localização do curso6.
O desempenho dos estudantes dentro de cada Área/Subárea foi classificado em quartos ( 1/4 ) para os concluintes e para os ingressantes em separado ( ver Figura 1 para o esquema de representação gráfica destas grandezas ).
As notas finais foram colocadas em ordem ascendente e no primeiro quarto foram colocadas as 25 \% notas mais baixas, limitadas superiormente pelo percentil 25 que vem a ser o primeiro quartil.
No segundo quarto, as 25 \% seguintes, limitadas superiormente pelo percentil 50 que vem a ser a mediana.
No terceiro quarto, as 25 \% notas após as encontradas no segundo quarto, ou seja, as primeiras 25 \% acima da mediana, limitadas superiormente pelo percentil 75 que vem a ser o terceiro quartil.
No último quarto foram alocadas as 25 \% maiores notas.
Potencialmente, poderia existir uma relação entre a nota do estudante e a sua percepção sobre a prova.
Calculando as médias das respostas segundo os quartos de notas, seria possível reconhecer a existência ( ou não ) de algum gradiente nas respostas, por exemplo, uma melhor nota correspondendo a uma percepção menos negativa.
Este gradiente seria traduzido, por exemplo, no gráfico 4.1a, pela diminuição das áreas \textbf{laranja} e vermelha ( Difícil e Muito difícil ) com o aumento do desempenho ( indo do primeiro ao último quarto ). 6 Curso, neste relatório, refere -se à unidade de análise para o Conceito ENADE e é caracterizado pela combinação de IES, Área/Subárea ( quando pertinente ) e município de habilitação.
A medida tradicional de “ grau de discriminação ” comparou os resultados dos estudantes do primeiro quarto ( 25 \% piores estudantes ) aos do último quarto ( 25 \% melhores ).
Esta comparação, no gráfico disponibilizado, foi idêntica à comparação da primeira barra ( primeiro quarto ) com a última barra ( último quarto ) de cada grupo ingressante ou concluinte nos gráficos “ segundo quartos de nota final ” ( gráficos com números ímpares neste capítulo ).
Por outro lado, a disponibilização das informações segundo os quartos de nota permitiu verificar o comportamento de dois outros grupos intermediários que na análise somente dos extremos tem que ser suposto monotônico.
A seguir, serão apresentados os principais resultados relativos aos nove itens avaliados. 4.1 Grau de dificuldade da prova 4.1.1 Formação Geral Os Gráficos 4.1a, 4.1b e 4.1c apresentam a distribuição das respostas à questão “ Qual o grau de dificuldade desta prova na parte de Formação Geral? ”, segundo quartos de nota final e grupo de estudantes concluintes e ingressantes que participaram do ENADE/2009 na área de TECNOLOGIA EM DESIGN DE MODA.
No total de concluintes, 17,0 \% declararam que a prova na parte da Formação Geral foi Difícil/Muito difícil, contra um valor acima para os ingressantes, 20,0 \%.
Tanto entre os concluintes, como entre os ingressantes, a proporção de indivíduos que classificaram esta parte como Difícil/Muito difícil diminui com o desempenho do estudante.
Por outro lado, a proporção que declarou a parte Fácil/Muito fácil foi ligeiramente maior entre os concluintes ( 17,5 \% ) do que entre os ingressantes ( 16,0 \% ).
Observou -se um crescimento desta proporção com o aumento no desempenho, tanto entre os concluintes, como entre os ingressantes.
Entre os concluintes, 65,5 \% classificaram esta parte com um nível Médio de dificuldade.
A proporção entre os ingressantes foi ligeiramente menor, 64,1 \%.
Os Gráficos 4.2a e 4.2b apresentam a distribuição das respostas à questão “ Qual o grau de dificuldade desta prova na parte de Formação Geral? ”, segundo a Região onde o curso está situado para os estudantes concluintes e ingressantes que participaram do ENADE/2009 na área de TECNOLOGIA EM DESIGN DE MODA.
Entre os concluintes, a maior proporção de estudantes que considerou a parte de Formação Geral da prova como de dificuldade média entre todas as Regiões brasileiras, foram os da Região Norte com 83,3 \% e entre os ingressantes na Nordeste com 67,7 \%.
A Região com as menores proporções desta resposta foram respectivamente para os concluintes e ingressantes, a Centro-Oeste ( 61,5 \% ) e a Norte ( 60,0 \% ). 4.1.2 Componente Específico Os Gráficos 4.3a, 4.3b e 4.3c apresentam a distribuição das respostas à questão “ Qual o grau de dificuldade desta prova na parte do Componente Específico? ”, segundo quartos de nota final e grupo de estudantes concluintes e ingressantes que participaram do ENADE/2009 na área de TECNOLOGIA EM DESIGN DE MODA.
Em geral, para as desagregações consideradas, os estudantes classificaram o Componente Específico como mais difícil que a Formação Geral.
No total de concluintes, 14,6 \% declararam que a prova na parte do Componente Específico foi Difícil/Muito difícil, sendo um valor muito maior para os ingressantes, 31,3 \%.
Entre os concluintes, a proporção de indivíduos que classificaram esta parte como Difícil/Muito difícil diminuiu com o desempenho do estudante.
Entre os concluintes 1,4 \% considerou esta parte Muito fácil.
Entre os ingressantes 0,5 \% também o fizeram.
A proporção que declarou a parte Fácil/Muito fácil foi maior entre os concluintes ( 19,9 \% ) do que entre os ingressantes ( 9,8 \% ).
Entre os concluintes a proporção dos que consideraram esta parte como fácil cresce com o desempenho.
Entre os concluintes, 65,5 \% classificaram esta parte com um nível Médio de dificuldade.
A proporção entre os ingressantes foi ligeiramente menor, 59,0 \%.
Entre os concluintes, esta proporção aumenta com o desempenho, indo de 58,0 \% no primeiro quarto de notas a 64,6 \% no último quarto.
Os Gráficos 4.4a e 4.4b apresentam a distribuição das respostas à questão “ Qual o grau de dificuldade desta prova na parte do Componente Específico? ”, segundo a Região onde o curso está situado para os estudantes concluintes e ingressantes que participaram do ENADE/2009 na área de TECNOLOGIA EM DESIGN DE MODA.
Considerando -se as respostas dadas pelos estudantes concluintes de cada Região, observa -se que somente nas Regiões Nordeste ( 12,0 \% ) e Centro-Oeste ( 10,3 \% ) o percentual de estudantes que avaliaram a parte do Componente Específico da prova como Difícil ou Muito difícil foi inferior à média nacional ( 14,6 \% ).
As Regiões Norte e Nordeste foram as únicas a não terem a declaração de Muito difícil.
Entre os concluintes apenas 1,4 \% declararam esta parte como Muito fácil e entre os ingressantes apenas 0,5 \% o fizeram.
Nenhum estudante da Região Norte classificou essa parte como Muito fácil.
Entre os concluintes, a maior proporção de estudantes que considerou a parte do Componente Específico da prova como de dificuldade média entre todas as Regiões brasileiras, foi na Região Norte com 83,3 \% e entre os ingressantes também com 66,7 \%.
A Região com as menores proporções desta resposta foram respectivamente para os concluintes e ingressantes, a Sul ( 63,3 \% ) e a Sudeste ( 57,2 \% ). 4.2 Extensão da prova em relação ao tempo total Os Gráficos 4.5a, 4.5b e 4.5c apresentam a distribuição das respostas à questão sobre a extensão da prova em relação ao tempo total, segundo quartos de nota final e grupo de estudantes concluintes e ingressantes que participaram do ENADE/2009 na área de TECNOLOGIA EM DESIGN DE MODA.
Tanto para concluintes quanto para ingressantes, a maioria considerou a extensão da prova como Adequada em relação ao tempo destinado à sua resolução, respectivamente 67,8 \% e 71,3 \%.
Entre os concluintes esta proporção cresceu com o desempenho.
A proporção dos que consideraram a extensão da prova Muito curta foi ligeiramente menor entre os concluintes ( 0,3 \% ) do que entre os ingressantes ( 0,4 \% ).
A proporção dos que consideraram a extensão da prova Curta/Muito curta vis-à-vis ao tempo total foi menor entre os concluintes ( 3,8 \% ) do que entre os ingressantes ( 4,5 \% ).
Entre os concluintes, somente estudantes com desempenho no primeiro, terceiro e no último quartos declararam que a extensão da prova teria sido Muito curta com relação ao tempo total.
Entre os ingressantes, esta declaração só não aconteceu no último quarto.
A maior proporção de concluintes declarando que a prova foi Curta/Muito curta foi entre os estudantes do terceiro quarto, 4,9 \% e entre os ingressantes entre os do segundo quarto, 4,7 \%.
A proporção dos que consideraram a extensão da prova Muito longa foi maior entre os concluintes ( 7,9 \% ) do que entre os ingressantes ( 6,7 \% ).
A proporção dos que consideraram a extensão da prova Longa foi menor entre os concluintes ( 20,4 \% ) do que entre os ingressantes ( 17,6 \% ).
Os Gráficos 4.6a e 4.6b apresentam a distribuição das respostas à questão sobre a extensão da prova em relação ao tempo total, segundo a Região onde o curso está situado para os estudantes concluintes e ingressantes que participaram do ENADE/2009 na área de TECNOLOGIA EM DESIGN DE MODA.
Em relação à percepção sobre o tempo da prova ser adequado, somente os concluintes da Região Sul ( 76,1 \% ) se situaram acima da média ( 67,8 \% ) para esta resposta, enquanto que entre os ingressantes, os da Região Nordeste ( 64,0 \% ) e Centro-Oeste ( 66,8 \% ) ficaram abaixo da média ( 71,3 \% ).
No outro extremo da classificação, com respeito à extensão da prova ser considerada como Curta/Muito curta, somente a Região Norte ( 16,7 \% ) e Centro-Oeste ( 7,8 \% ) entre os concluintes e a Região Centro-Oeste ( 10,4 \% ) entre os ingressantes apresentaram proporções acima da média nacional ( respectivamente 3,9 \% e 4,5 \% para concluintes e ingressantes ). 4.3 Compreensão dos Enunciados das Questões 4.3.1 Formação Geral Os Gráficos 4.7a, 4.7b e 4.7c apresentam a distribuição das respostas à questão “ Os enunciados das questões da prova na parte de Formação Geral, estavam claros e objetivos? ”, segundo quartos de nota final e grupo de estudantes concluintes e ingressantes que participaram do ENADE/2009 na área de TECNOLOGIA EM DESIGN DE MODA.
Com relação aos enunciados das questões da parte de Formação Geral, as opiniões foram positivas : a grande maioria dos concluintes ( 70,1 \% ) e ingressantes ( 72,9 \% ) consideraram claros e objetivos os enunciados de todas ou da maioria das questões.
Entre os concluintes, esta proporção cresce com o desempenho : 61,0 \% para o grupo com o pior desempenho e 79,0 \% para o de melhor.
Entre os ingressantes, esta proporção cresce com o mesmo comportamento : 65,8 \% para o grupo com o pior desempenho e 79,6 \% para o de melhor.
Entre ingressantes a combinação de “ Poucos ” e “ Não, nenhum ” decresceu com o desempenho.
Os Gráficos 4.8a e 4.8b apresentam a distribuição das respostas à questão “ Os enunciados das questões da prova na parte de Formação Geral, estavam claros e objetivos? ” segundo a Região onde o curso está situado para os estudantes concluintes e ingressantes que participaram do ENADE/2009 na área de TECNOLOGIA EM DESIG DE MODA.
Em Em termos regionais, as proporções de estudantes que consideraram claros e objetivos os enunciados de todas ou da maioria das questões variaram entre os concluintes, de 67,6 \% no Nordeste a 83,4 \% no Norte.
Já entre os ingressantes estas proporções variaram de 60,0 \% no Norte a 79,3 \% no Centro-Oeste.
Com relação às respostas “ Poucos ” e “ Não, nenhum ” as maiores proporções ocorreram para concluintes na Região Norte ( 16,7 \% ) e para ingressantes na Norte ( 20,0 \% ) 4.3.2 Componente Específico Os Gráficos 4.9a, 4.9b e 4.9c apresentam a distribuição das respostas à questão “ Os enunciados das questões da prova na parte de formação específica estavam claros e objetivos? ”, segundo quartos de nota final e grupo de estudantes concluintes e ingressantes que participaram do ENADE/2009 na área de TECNOLOGIA EM DESIGN DE MODA.
Com relação aos enunciados das questões da parte do Componente Específico, as opiniões foram também positivas : a grande maioria dos concluintes ( 73,8 \% ) e ingressantes ( 76,7 \% ) consideraram claros e objetivos os enunciados de todas ou da maioria das questões.
A proporção de ingressantes declarando que todos ou a maioria dos enunciados estavam claros e objetivos foi crescente com o desempenho indo de 69,2 \% para o primeiro quarto para 80,7 \% para os com melhor desempenho.
O número de estudantes que declararam que “ Não, nenhum ” dos enunciados das questões da parte do Componente Específico era claro e objetivo, foi um pouco menor do que na parte de Formação Geral : 0,7 \% entre os concluintes e 0,9 \% entre os ingressantes.
Os que responderam “ Poucos ”, são 7,2 \% entre os concluintes e 7,3 \% entre os ingressantes.
Os Gráficos 4.10a e 4.10b apresentam a distribuição das respostas à questão “ Os enunciados das questões da prova na parte de formação específica estavam claros e objetivos? ”, segundo a Região onde o curso está situado para os estudantes concluintes e ingressantes que participaram do ENADE/2009 na área de TECNOLOGIA DESIGN DE MODA.
Em termos regionais, as proporções de estudantes que consideraram claros e objetivos os enunciados de todas ou da maioria das questões variaram entre os concluintes, de 70,0 \% no Nordeste a 83,3 \% no Norte.
Já entre os ingressantes estas proporções variaram de 75,0 \% na Região Sudeste a 78,6 \% na Norte.
Com relação às respostas “ Poucos ” e “ Não, nenhum ” para concluintes a maior proporção ocorreu na Região Centro-Oeste com 9,5 \%.
Para os ingressantes a maior proporção ocorreu na Região Nordeste com 9,7 \%.
Tanto para concluintes como para ingressantes a opção que apresentou o maior percentual de escolha foi “ Sim, a maioria ” com 56,8 \% e 57,4 \% respectivamente. 4.4 Suficiência das informações/instruções fornecidas Os Gráficos 4.11a, 4.11b e 4.11c apresentam a distribuição das respostas à questão “ As informações/instruções fornecidas para a resolução das questões foram suficientes para resolvê -las? ”, segundo quartos de nota final e grupo de estudantes concluintes e ingressantes que participaram do ENADE/2009 na área de TECNOLOGIA EM DESIGN DE MODA.
Quando perguntados sobre as informações/instruções fornecidas para a resolução das questões, 80,4 \% dos ingressantes e 84,1 \% dos concluintes de todo o Brasil responderam que estas eram ou excessivas ou suficientes em todas ou na maioria das questões, demonstrando uma avaliação positiva da prova.
Levando -se em conta o desempenho dos participantes, percebeu -se consideráveis diferenças entre as opiniões de estudantes, ingressantes e concluintes.
A proporção de estudantes com a opinião supracitada foi crescente com o desempenho, tanto para concluintes como para ingressantes.
Entre os concluintes a proporção vai de 78,5 \% para os de pior desempenho a 88,7 \% para os de melhor desempenho.
Entre os ingressantes, as proporções equivalentes são, respectivamente, 75,7 \% a 87,2 \%.
O complemento desta opinião, ou seja, os estudantes que responderam “ Somente em algumas ” ou “ Em nenhuma delas ”, apresentou uma proporção decrescente com o desempenho tanto para os concluintes como para os ingressantes.
A resposta “ Não, em nenhuma delas ” entre os ingressantes, não apareceu no grupo de melhor desempenho.
Os Gráficos 4.12a e 4.12b apresentam a distribuição das respostas à questão “ As informações/instruções fornecidas para a resolução das questões foram suficientes para resolvê -las? ”, segundo a Região onde o curso está situado para os estudantes concluintes e ingressantes que participaram do ENADE/2009 na área de TECNOLOGIA EM DESIGN DE MODA.
A Região Sudeste foi a que apresentou o percentual mais baixo, com 82,8 \% entre os concluintes emitindo tal opinião.
E a Região Norte foi a que apresentou o percentual mais baixo, com 73,3 \% dos ingressantes emitindo tal opinião. 4.5 Dificuldade encontrada ao responder à prova Os Gráficos 4.13a, 4.13b e 4.13c apresentam a distribuição das respostas à questão “ Você se deparou com alguma dificuldade ao responder à prova? ”, segundo quartos de nota final e grupo de estudantes concluintes e ingressantes que participaram do ENADE/2009 na área de TECNOLOGIA EM DESIGN DE MODA.
Perguntados sobre as dificuldades com as quais se depararam ao responder à prova, 48,6 \% dos ingressantes apontaram como principal dificuldade o “ Desconhecimento do conteúdo ”, já os concluintes apontaram como principal dificuldade a “ Forma diferente de abordagem do conteúdo ”, 41,4 \%.
Na verdade para a quase totalidade das áreas, ingressantes apresentaram uma proporção de estudantes bem maior do que entre os concluintes declarando como principal dificuldade o desconhecimento da matéria.
O desconhecimento do conteúdo foi apontado somente por 14,1 \% dos concluintes, valores sem muita variação entre os três primeiros grupos de desempenho.
A forma diferente de abordagem do conteúdo foi considerada como maior dificuldade por 27,3 \% dos ingressantes, enquanto que a falta de motivação foi a principal dificuldade para 10,9 \% dos ingressantes.
Em termos nacionais, a proporção de concluintes e ingressantes que declararam não ter tido qualquer dificuldade para responder à prova foi, respectivamente, de 15,6 \% e 8,5 \%.
Os Gráficos 4.14a e 4.14b apresentam a distribuição das respostas à questão “ Você se deparou com alguma dificuldade ao responder à prova? ”, segundo a Região onde o curso está situado para os estudantes concluintes e ingressantes que participaram do ENADE/2009 na área de TECNOLOGIA EM DESIGN DE MODA.
Em relação às brasileiras, o percentual de ingressantes que apontaram o desconhecimento do conteúdo como principal dificuldade ao responder à prova variou de 13,3 \% na Região Norte, a 56,4 \% na Sul.
Entre os concluintes, tal percentual variou de 11,4 \% na Região Centro-Oeste a 33,3 \%, na Norte.
A forma diferente de abordagem de conteúdo foi citada como maior dificuldade por 16,7 \% dos concluintes na Região Norte chegando a 44,5 \% na Região Sul.
Já entre os ingressantes, os valores foram de 23,9 \% na Região Sul a 40,0 \% na Região Norte.
O percentual de concluintes que citaram a falta de motivação como principal dificuldade variou de 12,8 \% na Região Sul a 28,5 \% na Região Centro-Oeste.
Entre os ingressantes, as estatísticas correspondentes foram de 6,7 \% na Região Norte a 13,0 \%, na Região Centro-Oeste.
Os que declararam que não tiveram qualquer tipo de dificuldade variaram, entre os concluintes, de 0,0 \% na Região Norte a 17,2 \% na Região Sudeste.
Entre os ingressantes este percentual variou de 6,9 \% na Região Sudeste a 40,0 \% na Região Norte. 4.6 Percepção sobre os conteúdos das questões objetivas da prova Os Gráficos 4.15a, 4.15b e 4.15c apresentam a distribuição das respostas à questão “ Considerando apenas as questões objetivas da prova, você percebeu que :... ”, segundo quartos de nota final e grupo de estudantes concluintes e ingressantes que participaram do ENADE/2009 na área de TECNOLOGIA EM DESIGN DE MODA.
As respostas possíveis para esta questão estavam relacionadas à cobertura do conteúdo das questões objetivas durante o curso.
A proporção de ingressantes ( 54,0 \% ) que declarou que não estudou ainda a maioria desses conteúdos foi bem maior do que entre os concluintes ( 6,4 \% ).
Esta proporção foi decrescente tanto para os concluintes como para os ingressantes como função do desempenho.
Por outro lado, a proporção de estudantes que declarou que estudou e aprendeu todos ou a maioria desses conteúdos, entre os ingressantes foi de 30,0 \% e entre os concluintes 71,8 \%.
Esta proporção foi crescente como função do desempenho, tanto para os concluintes como para os ingressantes.
A proporção de estudantes que declarou ter estudado a maioria desses conteúdos, mas não os ter aprendido foi bem maior entre os concluintes ( 10,4 \% ) do que entre os ingressantes ( 4,9 \% ).
Os Gráficos 4.16a e 4.16b apresentam a distribuição das respostas à questão “ Considerando apenas as questões objetivas da prova, você percebeu que :... ”, para os estudantes concluintes e ingressantes que participaram do ENADE/2009 na área de TECNOLOGIA EM DESIGN DE MODA, segundo a Região onde o curso está situado.
Em relação às Regiões brasileiras, os concluintes que afirmaram não ter estudado ainda a maioria dos conteúdos variaram de 4,9 \% na Região Nordeste a 66,7 \% na Norte.
Entre os ingressantes, os percentuais variaram de 35,7 \% na Região Norte a 62,8 \% na Sul.
Em todas as Regiões, exceto a Região Norte, a maioria dos concluintes afirmou ter estudado e aprendido todos ou muitos dos conteúdos, com percentuais que variaram de 33,3 \% na Região Norte a 77,6 \% na Centro-Oeste.
Para os ingressantes tal percentual variou de 24,9 \% na Região Sul a 39,8 \% na Nordeste. 4.7 Tempo gasto para concluir a prova Os Gráficos 4.17a, 4.17b e 4.17c apresentam a distribuição das respostas à questão “ Qual foi o tempo gasto para concluir a prova? ”, segundo quartos de nota final e grupo de estudantes concluintes e ingressantes que participaram do ENADE/2009 na área de TECNOLOGIA EM DESIGN DE MODA.
Quando perguntados quanto ao tempo de conclusão da prova, o grupo modal respondeu entre uma e duas horas, tanto entre os concluintes ( 47,3 \% ) quanto entre os ingressantes ( 53,9 \% ).
A segunda maior incidência de respostas foi entre duas e três horas, igualmente para ambos os grupos : concluintes ( 32,4 \% ) e ingressantes ( 26,8 \% ).
A proporção daqueles que declararam ter usado as quatro horas, mas não terem conseguido terminar foi bem baixa : 0,7 \% entre os ingressantes e 1,2 \% entre os concluintes.
Entre os concluintes, não se percebeu uma tendência monotônica do tempo necessário para concluir a prova como função do desempenho.
Ao contrário, entre os ingressantes, notou -se uma tendência monotônica com o desempenho : a proporção de estudantes declarando menos de 2 horas é decrescente com o desempenho, caindo de 79,8 \% no primeiro grupo de notas a 56,5 \% no último grupo.
Os Gráficos 4.18a e 4.18b apresentam a distribuição das respostas à questão “ Qual foi o tempo gasto para concluir a prova? ”, segundo a Região onde o curso está situado para os estudantes concluintes e ingressantes que participaram do ENADE/2009 na área de TECNOLOGIA EM DESIGN DE MODA.
A proporção de concluintes que gastaram entre uma e duas horas foi maior na Região Norte ( 66,7 \% ) e menor na Região Centro-Oeste ( 45,5 \% ).
Para ingressantes, o maior percentual para este tempo ocorreu na Região Sul ( 57,1 \% ) e a menor na Região Centro-Oeste ( 46,6 \% ).
Na Região Norte nenhum estudante declarou mais de 3 horas.
Na Região Sul nenhum estudante declarou ter usado as 4 horas e não ter conseguido terminar a prova.

% matched lemmas: laranja
\end{textsample}
